
\begin{table*}[htbp]
\caption{Experimental Results of 4 Small LLMs Evaluated across 8 Benchmarks}
\label{tab3.2.1_llama}
\centering
\begin{tabular}{c | c | cccccccc | c} 
\hline
\hline

\rule{0pt}{2.2ex} Model & A-W Quant Type & ARC-C & ARC-E & BoolQ & HellaS & LamOp & Piqa & WinoG & MMLU & Mean \\

\hline 

\rule{0pt}{2.2ex} \multirow{10}{*}{\rotatebox{90}{Llama2-7B}} &  BF16 & 45.65 & 74.41 & 77.74 & 75.99 & 73.67 & 79.11 & 69.06 & 46.52 & 67.77 \\

\cline{2-11}
\rule{0pt}{2.2ex} & NVFP4 & 44.54 & 73.32 & 76.02 & 75.02 & 73.69 & 77.37 & 70.48 & 41.51 & 66.49\\

& \cellcolor{mygray}{\color{gray}\qquad{\small --- Acc Drop}} & \cellcolor{mygray}{\color{red}-1.11} & \cellcolor{mygray}{\color{red}-1.09} & \cellcolor{mygray}{\color{red}-1.72} & \cellcolor{mygray}{\color{red}-0.97} & \cellcolor{mygray}{\color{blue}+0.02} & \cellcolor{mygray}{\color{red}-1.74} & \cellcolor{mygray}{\color{blue}+1.42} & \cellcolor{mygray}{\color{red}-5.01} & \cellcolor{mygray}{\color{red}-1.28} \\

\rule{0pt}{2.2ex} & NVFP4+PTS & 44.62 & 72.01 & 76.67 & 74.92 & 73.16 & 77.91 & 66.69 & 43.43 & 66.18 \\

& \cellcolor{mygray}{\color{gray}\qquad{\small --- Acc Drop}} & \cellcolor{mygray}{\color{red}-1.03} & \cellcolor{mygray}{\color{red}-2.40} & \cellcolor{mygray}{\color{red}-1.07} & \cellcolor{mygray}{\color{red}-1.07} & \cellcolor{mygray}{\color{red}-0.51} & \cellcolor{mygray}{\color{red}-1.20} & \cellcolor{mygray}{\color{red}-2.37} & \cellcolor{mygray}{\color{red}-3.09} & \cellcolor{mygray}{\color{red}-1.59} \\

\rule{0pt}{2.2ex} & HiF4 & 44.63 & 74.08 & 76.73 & 75.00 & 72.65 & 78.40 & 68.39 & 44.54 & 66.80 \\

& \cellcolor{mygray}{\color{gray}\qquad{\small --- Acc Drop}} & \cellcolor{mygray}{\color{red}-1.02} & \cellcolor{mygray}{\color{red}-0.33} & \cellcolor{mygray}{\color{red}-1.01} & \cellcolor{mygray}{\color{red}-0.99} & \cellcolor{mygray}{\color{red}-1.02} & \cellcolor{mygray}{\color{red}-0.71} & \cellcolor{mygray}{\color{red}-0.67} & \cellcolor{mygray}{\color{red}-1.98} & \cellcolor{mygray}{\color{red}-0.97}\\

\rule{0pt}{2.2ex} & HiF4+HiGPTQ & 45.31 & 73.61 & 76.58 & 74.91 & 73.66 & 77.86 & 68.99 & 43.89 & 66.85 \\

& \cellcolor{mygray}{\color{gray}\qquad{\small --- Acc Drop}} & \cellcolor{mygray}{\color{red}-0.34} & \cellcolor{mygray}{\color{red}-0.80} & \cellcolor{mygray}{\color{red}-1.16} & \cellcolor{mygray}{\color{red}-1.08} & \cellcolor{mygray}{\color{red}-0.01} & \cellcolor{mygray}{\color{red}-1.25} & \cellcolor{mygray}{\color{red}-0.07} & \cellcolor{mygray}{\color{red}-2.63} & \cellcolor{mygray}{\color{red}-0.92}\\

\hline 

\rule{0pt}{2.2ex} \multirow{10}{*}{\rotatebox{90}{LLama3-8B}} & BF16 & 53.41 & 77.78 & 81.16 & 79.15 & 75.65 & 80.85 & 72.93 & 66.55 & 73.44 \\

\cline{2-11}
\rule{0pt}{2.2ex} & NVFP4 & 48.04 & 73.74 & 78.35 & 76.56 & 74.21 & 78.56 & 71.43 & 61.64 & 70.32 \\

& \cellcolor{mygray}{\color{gray}\qquad{\small --- Acc Drop}} & \cellcolor{mygray}{\color{red}-5.37} & \cellcolor{mygray}{\color{red}-4.04} & \cellcolor{mygray}{\color{red}-2.81} & \cellcolor{mygray}{\color{red}-2.59} & \cellcolor{mygray}{\color{red}-1.44} & \cellcolor{mygray}{\color{red}-2.29} & \cellcolor{mygray}{\color{red}-1.50} & \cellcolor{mygray}{\color{red}-4.91} & \cellcolor{mygray}{\color{red}-3.12} \\

\rule{0pt}{2.2ex} & NVFP4+PTS & 49.74 & 74.54 & 77.46 & 77.60 & 74.85 & 79.76 & 70.80 & 62.81 & 70.95\\

& \cellcolor{mygray}{\color{gray}\qquad{\small --- Acc Drop}} & \cellcolor{mygray}{\color{red}-3.67} & \cellcolor{mygray}{\color{red}-3.24} & \cellcolor{mygray}{\color{red}-3.70} & \cellcolor{mygray}{\color{red}-1.55} & \cellcolor{mygray}{\color{red}-0.80} & \cellcolor{mygray}{\color{red}-1.09} & \cellcolor{mygray}{\color{red}-2.13} & \cellcolor{mygray}{\color{red}-3.74} & \cellcolor{mygray}{\color{red}-2.49} \\

\rule{0pt}{2.2ex} & HiF4 & 51.15 & 76.73 & 79.85 & 77.84 & 73.76 & 79.09 & 71.79 & 63.37 & 71.70 \\

& \cellcolor{mygray}{\color{gray}\qquad{\small --- Acc Drop}} &  \cellcolor{mygray}{\color{red}-2.26} & \cellcolor{mygray}{\color{red}-1.05} & \cellcolor{mygray}{\color{red}-1.31} & \cellcolor{mygray}{\color{red}-1.31} & \cellcolor{mygray}{\color{red}-1.89} & \cellcolor{mygray}{\color{red}-1.76} & \cellcolor{mygray}{\color{red}-1.14} & \cellcolor{mygray}{\color{red}-3.18} & \cellcolor{mygray}{\color{red}-1.74}\\

\rule{0pt}{2.2ex} & HiF4+HiGPTQ & 52.18 & 77.46 & 80.40 & 77.47 & 72.85 & 79.60 & 71.87 & 63.55 & 71.92 \\

& \cellcolor{mygray}{\color{gray}\qquad{\small --- Acc Drop}} &  \cellcolor{mygray}{\color{red}-1.23} & \cellcolor{mygray}{\color{red}-0.32} & \cellcolor{mygray}{\color{red}-0.76} & \cellcolor{mygray}{\color{red}-1.68} & \cellcolor{mygray}{\color{red}-2.80} & \cellcolor{mygray}{\color{red}-1.25} & \cellcolor{mygray}{\color{red}-1.06} & \cellcolor{mygray}{\color{red}-3.00} & \cellcolor{mygray}{\color{red}-1.52}\\

\hline 

\rule{0pt}{2.2ex} \multirow{10}{*}{\rotatebox{90}{Qwen2.5-14B}} &  BF16 & 58.96 & 79.34 & 85.54 & 82.94 & 74.31 & 81.88 & 74.74 & 80.17 & 77.24\\

\cline{2-11}
\rule{0pt}{2.2ex} & NVFP4 & 58.53 & 80.56 & 83.49 & 81.38 & 72.95 & 81.72 & 74.43 & 76.53 & 76.20\\

& \cellcolor{mygray}{\color{gray}\qquad{\small --- Acc Drop}} & \cellcolor{mygray}{\color{red}-0.43} & \cellcolor{mygray}{\color{blue}+1.22} & \cellcolor{mygray}{\color{red}-2.05} & \cellcolor{mygray}{\color{red}-1.56} & \cellcolor{mygray}{\color{red}-1.36} & \cellcolor{mygray}{\color{red}-0.16} & \cellcolor{mygray}{\color{red}-0.31} & \cellcolor{mygray}{\color{red}-3.64} & \cellcolor{mygray}{\color{red}-1.04}\\

\rule{0pt}{2.2ex} & NVFP4+PTS & 56.57 & 80.18 & 85.72 & 81.45 & 72.87 & 81.66 & 73.16 & 78.63 & 76.28\\

& \cellcolor{mygray}{\color{gray}\qquad{\small --- Acc Drop}} & \cellcolor{mygray}{\color{red}-2.39} & \cellcolor{mygray}{\color{blue}+0.84} & \cellcolor{mygray}{\color{blue}+0.18} & \cellcolor{mygray}{\color{red}-1.49} & \cellcolor{mygray}{\color{red}-1.44} & \cellcolor{mygray}{\color{red}-0.22} & \cellcolor{mygray}{\color{red}-1.58} & \cellcolor{mygray}{\color{red}-1.54} & \cellcolor{mygray}{\color{red}-0.96}\\

\rule{0pt}{2.2ex} & HiF4 & 58.66 & 81.55 & 85.86 & 81.45 & 72.74 & 81.29 & 73.72 & 78.62 & 76.74 \\

& \cellcolor{mygray}{\color{gray}\qquad{\small --- Acc Drop}} & \cellcolor{mygray}{\color{red}-0.30} & \cellcolor{mygray}{\color{blue}+2.21} & \cellcolor{mygray}{\color{blue}+0.32} & \cellcolor{mygray}{\color{red}-1.49} & \cellcolor{mygray}{\color{red}-1.57} & \cellcolor{mygray}{\color{red}-0.59} & \cellcolor{mygray}{\color{red}-1.02} & \cellcolor{mygray}{\color{red}-1.55} & \cellcolor{mygray}{\color{red}-0.50}\\

\rule{0pt}{2.2ex} & HiF4+HiGPTQ & 60.71 & 83.04 & 86.27 & 81.38 & 73.02 & 81.67 & 75.06 & 78.65 & 77.48 \\

& \cellcolor{mygray}{\color{gray}\qquad{\small --- Acc Drop}} & \cellcolor{mygray}{\color{blue}+1.75} & \cellcolor{mygray}{\color{blue}+3.70} & \cellcolor{mygray}{\color{blue}+0.73} & \cellcolor{mygray}{\color{red}-1.56} & \cellcolor{mygray}{\color{red}-1.29} & \cellcolor{mygray}{\color{red}-0.21} & \cellcolor{mygray}{\color{blue}+0.32} & \cellcolor{mygray}{\color{red}-1.52} & \cellcolor{mygray}{\color{blue}+0.24}\\

\hline 

\rule{0pt}{2.2ex} \multirow{10}{*}{\rotatebox{90}{Mistral-7B}} &  BF16 &  52.39 & 78.37 & 82.17 & 80.50 & 75.14 & 82.21 & 74.11 & 63.30 & 73.52\\

\cline{2-11}
\rule{0pt}{2.2ex} & NVFP4 & 28.41 & 26.30 & 56.09 & 26.69 & 0.0 & 48.69 & 48.30 & 26.79 & 32.66\\

& \cellcolor{mygray}{\color{gray}\qquad{\small --- Acc Drop}} & \cellcolor{mygray}{\color{red}-23.98} & \cellcolor{mygray}{\color{red}-52.07} & \cellcolor{mygray}{\color{red}-26.08} & \cellcolor{mygray}{\color{red}-53.81} & \cellcolor{mygray}{\color{red}-75.14} & \cellcolor{mygray}{\color{red}-33.52} & \cellcolor{mygray}{\color{red}-25.81} & \cellcolor{mygray}{\color{red}-36.51} & \cellcolor{mygray}{\color{red}-40.87} \\

\rule{0pt}{2.2ex} & NVFP4+PTS & 50.77 & 77.19 & 80.86 & 79.71 & 73.96 & 80.85 & 72.45 & 61.15 & 72.12 \\

& \cellcolor{mygray}{\color{gray}\qquad{\small --- Acc Drop}} & \cellcolor{mygray}{\color{red}-1.62} & \cellcolor{mygray}{\color{red}-1.18} & \cellcolor{mygray}{\color{red}-1.31} & \cellcolor{mygray}{\color{red}-0.79} & \cellcolor{mygray}{\color{red}-1.18} & \cellcolor{mygray}{\color{red}-1.36} & \cellcolor{mygray}{\color{red}-1.66} & \cellcolor{mygray}{\color{red}-2.15} & \cellcolor{mygray}{\color{red}-1.41}  \\

\rule{0pt}{2.2ex} & HiF4 & 51.41 & 76.85 & 81.35 & 79.41 & 74.11 & 81.01 & 72.06 & 61.63 & 72.23 \\

& \cellcolor{mygray}{\color{gray}\qquad{\small --- Acc Drop}} & \cellcolor{mygray}{\color{red}-0.98} & \cellcolor{mygray}{\color{red}-1.52} & \cellcolor{mygray}{\color{red}-0.82} & \cellcolor{mygray}{\color{red}-1.09} & \cellcolor{mygray}{\color{red}-1.03} & \cellcolor{mygray}{\color{red}-1.20} & \cellcolor{mygray}{\color{red}-2.05} & \cellcolor{mygray}{\color{red}-1.67} & \cellcolor{mygray}{\color{red}-1.29}\\

\rule{0pt}{2.2ex} & HiF4+HiGPTQ & 51.84 & 77.74 & 82.22 & 79.48 & 73.31 & 81.91 & 73.40 & 61.53 & 72.68 \\

& \cellcolor{mygray}{\color{gray}\qquad{\small --- Acc Drop}} & \cellcolor{mygray}{\color{red}-0.55} & \cellcolor{mygray}{\color{red}-0.63} & \cellcolor{mygray}{\color{blue}+0.05} & \cellcolor{mygray}{\color{red}-1.02} & \cellcolor{mygray}{\color{red}-1.83} & \cellcolor{mygray}{\color{red}-0.30} & \cellcolor{mygray}{\color{red}-0.71} & \cellcolor{mygray}{\color{red}-1.77} & \cellcolor{mygray}{\color{red}-0.84}\\

\hline
\hline
\end{tabular}
\end{table*}

\section{Language Model Inference with HiFloat4}

This section conducts a comprehensive inference evaluation of 
Post-Training Quantization (PTQ) across diverse tasks and LLM 
architectures on HiF4 and NVFP4 formats.  

\subsection{Post-Training Quantization for LLMs}

To demonstrate the inherent precision of the 4-bit formats, 
we prioritize using HiF4 direct-cast, NVFP4 direct-cast, 
and the NVFP4 + PTS approach for LLM inference, as described 
in Section~\ref{Quantization Error and Dot Product}. 
Meanwhile, to show that many existing PTQ methods can be 
adapted to enhance HiF4 accuracy with minor changes, 
we developed a tailored PTQ method based on the vanilla 
GPTQ~\cite{Frantar2023}, termed HiGPTQ, 
to exploit the fine-grained structure of the HiF4 BFP format. 

\subsection{Experiments on Small LLMs}
We begin by evaluating our method on relatively smaller LLMs 
with the parameter size ranging from 7B to 14B. 
Detailed results and the corresponding inference accuracy degradation, 
denoted as Acc Drop, across eight benchmarks are presented in 
Table~\ref{tab3.2.1_llama}. 
To provide an overall view of quantization errors, 
Table~\ref{tab3.2.3_average} reports the average inference accuracy 
across all evaluated models (w/ and w/o the crashed Mistral-7B) for 
NVFP4, NVFP4+PTS, HiF4, and HiF4+HiGPTQ. 

\noindent\textbf{Models:} 
We evaluate 4-bit quantization accuracy across a diverse 
collection of LLM architectures and model scales, 
including LLaMA2-7B~\cite{Touvron2023}, 
LLaMA3-8B~\cite{Dubey2024}, 
Qwen2.5-14B~\cite{Qwen2025}, 
and Mistral-7B~\cite{Jiang2023}. 
These models cover a broad range of architectural designs, 
including Multi-Head Attention (MHA)~\cite{Vaswani2017} 
and Grouped-Query Attention (GQA)~\cite{Ainslie2023}, 
as well as multiple forms of Feed-Forward Networks 
(FFNs)~\cite{Shazeer2020}. 


\begin{table}[htbp]
\caption{Average Inference Accuracy for Small LLMs}
\label{tab3.2.3_average}
\centering
%\fontsize{7.5}{10}\selectfont
\begin{tabular}{c|c|cccc}
    
\hline
\hline

\rule{0pt}{2.2ex} \# LLM Models     & BF16  & NVFP4              & \makecell{NVFP4+\\PTS} & HiF4               & \makecell{HiF4+\\HiGPTQ} \\
\hline 

\rule{0pt}{2.2ex} 4 (w/ Mistral-7B) & 72.99 & 61.42              & 71.38                  & \textbf{71.87}     & 72.23 \\
\rule{0pt}{2.2ex} — Acc Drop        & -     & \color{red}{Crash} & \color{red}{-1.61}     & \color{red}{-1.12} & \color{red}{-0.76} \\
\hline

\rule{0pt}{2.2ex} 3(w/o Mistral-7B) & 72.82 & 71.01              & 71.14                  & \textbf{71.75}     & 72.08 \\
\rule{0pt}{2.2ex} — Acc Drop        & -     & \color{red}{-1.81} & \color{red}{-1.68}     & \color{red}{-1.07} & \color{red}{-0.74} \\
\hline

\hline
\end{tabular}
\end{table}

% (w/o Mistral-7B)
\begin{table*}[htbp]
\caption{Experimental Results of DeepSeek-V3.1 and LongCat Evaluated across 10 Benchmarks}
\label{tab3.2.1_deepseek}
\centering

\begin{tabular}{c | c | cccccccccc | c} 
\hline
\hline

\rule{0pt}{2.2ex} Model & A-W Quant Type & ARC-C & ARC-E & BoolQ & HellaS  & Piqa & WinoG &Gsm8K & MMLU  & Math500  & CMMLU & Mean \\

\hline 

\rule{0pt}{2.2ex} \multirow{7}{*}{\rotatebox{90}{\makecell{DeepSeek-V3.1\\671B}}} &  BF16 & 79.91 & 84.44 & 79.76 & 84.41 & 92.93 & 89.34 & 94.46 &  84.86 &75.00 & 89.28
    & 85.44 \\

\cline{2-13}
\rule{0pt}{2.2ex} & NVFP4 & 82.83 & 86.85 & 76.91 & 81.18 & 91.29 & 87.92 & 95.00 & 84.31 &74.00&87.76& 84.81\\

& \cellcolor{mygray}{\color{gray}\qquad{\small --- Acc Drop}} & \cellcolor{mygray}{\color{blue}+2.92} & \cellcolor{mygray}{\color{blue}+2.41
} & \cellcolor{mygray}{\color{red}-2.85} & \cellcolor{mygray}{\color{red}-3.23} & \cellcolor{mygray}{\color{red}-1.64} & \cellcolor{mygray}{\color{red}-1.42} & \cellcolor{mygray}{\color{blue}+0.54} & \cellcolor{mygray}{\color{red}-0.55} & \cellcolor{mygray}{\color{red}-1.00}& \cellcolor{mygray}{\color{red}-1.52} & \cellcolor{mygray}{\color{red}-0.63} \\



\rule{0pt}{2.2ex} & NVFP4+PTS & 81.55 & 87.06 & 77.83 & 81.00 & 91.29 & 88.00 & 95.30 & 
84.72 & 78.20 & 87.44
    & 85.24\\

& \cellcolor{mygray}{\color{gray}\qquad{\small --- Acc Drop}} & \cellcolor{mygray}{\color{blue}+1.64} & \cellcolor{mygray}{\color{blue}+2.62} & \cellcolor{mygray}{\color{red}-1.93} & \cellcolor{mygray}{\color{red}-3.41} & \cellcolor{mygray}{\color{red}-1.64} & \cellcolor{mygray}{\color{red}-1.34} & \cellcolor{mygray}{\color{blue}+0.84} & \cellcolor{mygray}{\color{red}-0.14} & \cellcolor{mygray}{\color{blue}+3.20} & \cellcolor{mygray}{\color{red}-1.84} & \cellcolor{mygray}{\color{red}-0.20} \\


\rule{0pt}{2.2ex} & HiF4 & 80.86 & 86.89& 78.87 & 86.60 & 92.44 & 89.11 & 95.75 & 84.77 & 80.20  & 88.71 & 86.42\\

& \cellcolor{mygray}{\color{gray}\qquad{\small --- Acc Drop}} & \cellcolor{mygray}{\color{blue}+0.95} & \cellcolor{mygray}{\color{blue}+2.45} & \cellcolor{mygray}{\color{red}-0.89} & \cellcolor{mygray}{\color{blue}+2.19} & \cellcolor{mygray}{\color{red}-0.49} & \cellcolor{mygray}{\color{red}-0.23} & \cellcolor{mygray}{\color{blue}+1.29} & \cellcolor{mygray}{\color{red}-0.09} & \cellcolor{mygray}{\color{blue}+5.20} & \cellcolor{mygray}{\color{red}-0.57} & \cellcolor{mygray}{\color{blue}+0.98}\\

\hline 
\hline

\hline 

\rule{0pt}{2.2ex} \multirow{7}{*}{\rotatebox{90}{\makecell{LongCat\\560B}}} &  BF16 & 84.38 & 86.64 & 66.85 & 82.09 & 91.46 & 80.27 & 95.91 & 59.19&84.80&81.65&81.32 \\

\cline{2-13}
\rule{0pt}{2.2ex} & NVFP4 &84.72&88.50&62.78&84.42&89.39&79.24&96.29&38.81&78.60&72.12&77.49
\\

& \cellcolor{mygray}{\color{gray}\qquad{\small --- Acc Drop}} & \cellcolor{mygray}{\color{blue}+0.34}  & \cellcolor{mygray}{\color{blue}+1.86}& \cellcolor{mygray}{\color{red}-4.07
} & \cellcolor{mygray}{\color{blue}+2.33} & \cellcolor{mygray}{\color{red}-2.07} & \cellcolor{mygray}{\color{red}-1.03} & \cellcolor{mygray}{\color{blue}+0.38} & \cellcolor{mygray}{\color{red}-20.38} & \cellcolor{mygray}{\color{red}-6.20}& \cellcolor{mygray}{\color{red}-9.53} & \cellcolor{mygray}{\color{red}-3.84} \\



\rule{0pt}{2.2ex} & NVFP4+PTS &87.04&89.26&63.55&84.17 &89.83&79.40&95.45&39.00&78.20&72.24 &77.81\\

& \cellcolor{mygray}{\color{gray}\qquad{\small --- Acc Drop}} & \cellcolor{mygray}{\color{blue}+2.66}  & \cellcolor{mygray}{\color{blue}+2.62}& \cellcolor{mygray}{\color{red}-3.30} & \cellcolor{mygray}{\color{blue}+2.08}  & \cellcolor{mygray}{\color{red}-1.63} & \cellcolor{mygray}{\color{red}-0.87} & \cellcolor{mygray}{\color{red}-0.46} & \cellcolor{mygray}{\color{red}-20.19} & \cellcolor{mygray}{\color{red}-6.60} & \cellcolor{mygray}{\color{red}-9.41} & \cellcolor{mygray}{\color{red}-3.51} \\


\rule{0pt}{2.2ex} & HiF4 &84.98
&87.95&71.19&84.04 &91.19&79.56&95.75&58.03&82.60&82.74&81.80\\

& \cellcolor{mygray}{\color{gray}\qquad{\small --- Acc Drop}} & \cellcolor{mygray}{\color{blue}+0.60}  & \cellcolor{mygray}{\color{blue}+1.31}& \cellcolor{mygray}{\color{blue}+4.34} & \cellcolor{mygray}{\color{blue}+1.95}  & \cellcolor{mygray}{\color{red}-0.27} & \cellcolor{mygray}{\color{red}-0.71} & \cellcolor{mygray}{\color{red}-0.16} & \cellcolor{mygray}{\color{red}-1.16} & \cellcolor{mygray}{\color{red}-2.20} & \cellcolor{mygray}{\color{blue}+1.09} & \cellcolor{mygray}{\color{blue}+0.48}\\

\hline 
\hline
\end{tabular}
\end{table*}


\noindent\textbf{Datasets:} 
We use a suite of widely adopted LLM benchmarks, 
including ARC-C(hallenge), ARC-E(asy)~\cite{Clark2018}, 
BoolQ~\cite{Clark2019}, HellaS(wag)~\cite{Zellers2019}, 
Lam(bada)Op(enAI)~\cite{Paperno2016}, 
Piqa~\cite{Bisk2020}, WinoG(rande)~\cite{Sakaguchi2019}, 
and MMLU~\cite{Hendrycks2021}. 
Together, these benchmarks span a broad spectrum of tasks, 
including commonsense and physical reasoning, 
cloze-style completion, pronoun/coreference resolution, 
fact verification, 
and multi-domain knowledge assessment. 

\noindent\textbf{Implementation Details:} 
Experiments in this section are all conducted using simulated 4-bit BFP 
formats based on Nvidia GPU and Huawei Ascend NPU. 
Specifically, before matrix multiplication, all linear layer tensors 
(with the exception of the embedding and LLM head layers) were converted 
from high-precision formats only to those that could be represented in 
HiF4 and NVFP4 formats. 
All experiments reported in this sub-section were conducted in parallel 
on two devices. 
For each device, we performed three runs using different random seeds. 
The final results are the average inference accuracy across the three runs 
and both devices to ensure that random variance is minimized. 
For a fair comparison, we evaluate all models and quantization settings 
within a unified framework, in which all experiments on a given benchmark 
employ the same prompt template for LLM inference. 

\noindent\textbf{Experimental Results and Analysis:} 
As shown in Tables~\ref{tab3.2.1_llama} and \ref{tab3.2.3_average}, 
HiF4 (direct-cast) consistently outperforms the prior state-of-the-art 
NVFP4 (direct-cast) and NVFP4+PTS across all four evaluated LLM models,  
demonstrating that the proposed HiF4 constitutes a more accurate 4-bit 
BFP format. 
We further observe that HiF4 is substantially more robust than NVFP4 
under direct conversion, 
attributable to its larger global and suitable local dynamic ranges. 
For instance, Mistral-7B exhibits a broader numerical distribution, 
under which direct-cast inference with NVFP4 results in model failure 
and accuracy degradation to near random-guess levels, 
\textit{i.e.}, inference crash. 
In contrast, HiF4 enables stable direct-cast inference under the same 
conditions. 
In addition, HiGPTQ further improves the inference 
accuracy of HiF4 across all four evaluated models. 
Notably, for the Qwen2.5-14B model, whose numerical distributions 
are optimized during training, 
the combination of HiF4 and HiGPTQ even surpasses the BF16 baseline. 

In summary, the proposed HiF4 format demonstrates higher inherent 
precision and stability than NVFP4. 
Many mature PTQ methods, such as GPTQ, can be adapted with minor 
modifications to further enhance the accuracy of HiF4. 


\subsection{Experiments on DeepSeek-V3.1 and LongCat} 

\noindent\textbf{Models:} 
We further evaluate 4-bit inference 
accuracy on larger LLMs, DeepSeek-V3.1 (671B)~\cite{Liu2024} 
and LongCat (560B)~\cite{Team2025}. 
These two models further encompass Multi-Head Latent Attention 
(MLA) and Mixture-of-Experts (MoE) architectures, 
which are commonly used in latest state-of-the-art LLMs. 

\noindent\textbf{Datasets:} 
In addition to the benchmarks mentioned above, 
our evaluation suite extends to several high-stakes domains, 
specifically Gsm8K~\cite{Cobbe2021}, 
Math500~\cite{Hendrycks2021a} 
and CMMLU~\cite{Li2024a} 
to ensure a holistic assessment of the model's capabilities. 

\noindent\textbf{Implementation Details:} 
We deployed the LongCat and DeepSeek-V3.1 models using the vLLM 
framework on clusters of 32 and 64 Ascend 910B NPUs, 
respectively, distributed across 4 and 8 nodes. 
To ensure a standardized accuracy assessment, 
the AISBench framework was employed for evaluation. 
All experiments were completed within a practical timeframe. 
For 4-bit evaluation, we converted the tensors in the MLA\_linear, 
MoE\_linear (excluding the gating network), and FFN\_linear layers 
to HiF4 and NVFP4 formats. 

\noindent\textbf{Experimental Results and Analysis:} 
Results in Table~\ref{tab3.2.1_deepseek} demonstrate that 
HiF4 (direct-cast) consistently surpasses the state-of-the-art 
NVFP4 format in 4-bit inference accuracy, 
including NVFP4 (direct-cast) and NVFP4+PTS. 
The superiority of HiF4 stems from its novel structure design 
as depicted in Figure~\ref{HiF4-Structure}, 
and more balanced features as outlined in Table~\ref{HiF4-NVFP4}, 
which provide greater robustness for deep learning. 
For instance, HiF4 successfully handles the relatively  
quantization-sensitive LLM LongCat, where NVFP4 fails and leads 
to accuracy crashes in some hard tasks.  
Furthermore, in specialized cases like the DeepSeek-V3.1 model, 
the inherent design of HiF4 allows it to rival the BF16 
baseline, highlighting its effectiveness as a high-fidelity 
4-bit BFP format. 

